\chapter{Future Work}\label{chap:futureWork}

While our heuristic is not comparable in effectiveness to the best current POCL heuristics for solvable and unsolvable instances, we have demonstrated that heuristics which account for all POCL search node information can be substantially more informed than heuristics currently used for POCL planning. Hence, avenues of future work to improve the runtime of our heuristic could lead to it outperforming other heuristics.\\


Firstly, the ILP model used for our heuristic is not efficient. By developing a model from scratch instead of using the model by~\cite{Imai2015ILPDeleteRelaxedPlanning} as a guideline, a more simple and effective ILP could be created. For example the $\mathit{init}$ and $\mathit{goal}$ steps do not have to be considered separately. Satisfaction of goals could be joined with the constraint requiring the satisfaction of preconditions, and the initial propositions would be established as effects of $\mathit{init}$, just like the effects of any other action. Action and step variables and constraints could also be combined, with an additional constraint requiring that all steps must be executed, as otherwise actions and steps are equivalent. In addition, action effect and step effect variables could be removed, as they are currently only used to ensure that propositions have an achiever, but this could be done by directly checking the positive effects of previously executed actions which occur after all possible deletions, in an inverse way to how preconditions are currently verified within our model. Finally, the way that an upper bound $ub$ is determined in our model is very relevant to runtime, as it lowers the number of values that time variables can hold, and tightens constraints on the objective value. Currently, the way that $ub$ is calculated assumes that the solution of the ILP will be a completely linear plan. However, our ILP model allows for actions to run in parallel if possible, and by accounting for this, a tighter upper bound on time variable values could be determined, improving the runtime of the ILP.\\


Each calculation of the ILP model for a new search node is done almost entirely independently in our implementation. The only information that is carried from previous calculations is the objective value, for use in determining an upper bound on the new objective value. However, the solution found by the ILP model at any search node will contain large amounts of important information about how search should progress from that node. For example, actions added to the plan by the ILP could be given higher priority for resolving flaws, as they are more likely to be required in a solution plan. The solution to the ILP of a parent search node could also sometimes be used to generate a feasible, but not necessarily optimal, solution for its children. This is a commonly used concept in linear programming called \textit{warmstarting}, which can significantly improve the runtime of ILP calculations.\\


There are also cases where evaluating our heuristic on a specific search node is not required. For example, in cases where a flaw has been repaired without a step being added, the only possible change in heuristic value from the parent node is if the new node has become unsolvable, which can often be determined by looking at the ILP solution of the parent node. By only using our more expensive heuristic for critical branching points in the search tree, for example adding a new step with delete effects to the plan, and using cheaper evaluations for less critical search nodes, performance could be improved.\\


The number of action and proposition variables created in our ILP could also be restricted by utilising information already obtained through the search process. For example, actions and propositions which can be found, using planning graphs or other means, to be either unreachable or not useful, do not require associated variables within the ILP. If there is knowledge of how many times a certain action might need to be executed or proposition might need to be achieved, this can also inform how many instance variables are created in the ILP. There might also be information about when a certain action or proposition can occur, for example if an action can at the earliest be executed at time step 4, then it's corresponding time variable in the ILP should have its domain restricted accordingly.\\


Different approaches could also be looked into for future work. For example, further relaxations, such as a linear relaxation of some version of our heuristic, could be promising. Alternatively, further restrictions could be applied, for example evaluating our heuristic for optimal planning, where the output plan must of minimal length. Given that our heuristic is admissible, it could easily be translated to optimal planning, and this could lead to more favourable comparisons with alternative heuristics. By changing the optimisation function of our ILP to minimise makespan instead of the number of actions, future works could also look into the use of our heuristic for minimum makespan planning.